% Заключение

В ходе выполнения курсовой работы были исследованы возможности Apache Spark для распределённой обработки, анализа и машинного обучения на больших данных. В качестве предметной области были выбраны сведения о дорожно-транспортных происшествиях в США за период с 2016 по 2023 год. Исходный набор US Accidents содержит 7\,728\,394 записи и 46 признаков, поэтому его обработка позволила рассмотреть основные этапы работы с крупными наборами данных в распределённой вычислительной среде.

На первом этапе исходные данные были загружены и обработаны с использованием Apache Spark и PySpark. После отбора необходимых полей и преобразования типов была сформирована таблица, сохранённая в формате Apache Iceberg. Количество записей после загрузки совпало с количеством строк исходного файла, что подтвердило отсутствие потерь данных на данном этапе. В ходе разведочного анализа были изучены распределения признаков, выявлены пропущенные значения, выбросы и аномальные погодные показатели. Для повышения качества данных применялись заполнение пропусков медианными значениями, проверка физических диапазонов, винзоризация экстремальных значений, объединение редких категорий и исключение признаков с высокой кардинальностью или низкой вариативностью. В результате был подготовлен набор данных, пригодный для последующего построения моделей машинного обучения.

На втором этапе средствами Apache Spark MLlib были рассмотрены две задачи машинного обучения. Для прогнозирования длительности происшествия была построена модель линейной регрессии. На тестовой выборке она показала RMSE 178,54 мин, MAE 85,00 мин и коэффициент детерминации $R^2 = 0,0472$. По сравнению с базовым предсказателем, возвращающим медиану обучающей выборки, среднеквадратичная ошибка уменьшилась на 4,85\,\%. Несмотря на то что модель превзошла базовый вариант, низкое значение $R^2$ указывает на ограниченную способность линейной зависимости описывать длительность происшествия. Следовательно, для повышения качества прогнозирования целесообразно исследовать нелинейные модели и способы преобразования целевой переменной.

В задаче классификации определялась принадлежность происшествия к категории тяжёлых: положительный класс соответствовал значениям \texttt{Severity} не ниже 3. Для решения задачи был использован алгоритм случайного леса. Из-за существенного дисбаланса исходных классов обучающая выборка балансировалась методом oversampling, в то время как тестовая выборка сохраняла естественное распределение данных. На тестовой выборке модель достигла значения AUC ROC 0,7970. При стандартном пороге классификации 0,5 полнота составила 0,7767. После изменения порога до 0,61 точность увеличилась до 0,4332 при полноте 0,6018, что позволило выполнить заданное ограничение на обнаружение тяжёлых происшествий. Полученные результаты показывают, что случайный лес способен выявлять зависимости, которые не описываются простой линейной моделью.

Сравнение результатов двух экспериментов показало, что выбор алгоритма должен определяться характером решаемой задачи и структурой данных. Линейная регрессия дала небольшой выигрыш относительно базового прогноза, тогда как случайный лес обеспечил более содержательное качество ранжирования объектов двух классов. При этом оценка важности признаков классификатора показала существенный вклад длительности и протяжённости происшествия. Эти характеристики становятся известны после возникновения события, поэтому построенная модель не может непосредственно рассматриваться как средство раннего предупреждения. Для такого применения необходимо изменить постановку задачи и исключить признаки, содержащие информацию о последствиях происшествия.

Таким образом, поставленная цель работы достигнута: рассмотрены основные этапы распределённой обработки и разведочного анализа набора US Accidents, выполнена подготовка данных и реализованы модели регрессии и бинарной классификации средствами Apache Spark MLlib. В ходе исследования были получены количественные оценки качества моделей и выявлены ограничения выбранных подходов. Практическая значимость выполненной работы заключается в демонстрации полного процесса анализа крупного набора данных -- от загрузки и очистки до обучения и оценки моделей -- с использованием распределённой вычислительной инфраструктуры.

В качестве дальнейшего развития исследования целесообразно рассмотреть применение более сложных алгоритмов, включая градиентный бустинг и ансамблевые методы, исследовать преобразование распределения целевой переменной в задаче регрессии, а также провести повторное обучение классификатора только на признаках, доступных непосредственно в момент возникновения происшествия. Дополнительно можно оценить влияние параметров распределённой обработки и объёма вычислительных ресурсов на время выполнения этапов подготовки данных и обучения моделей.